%! ~~~ Packages Setup ~~~ 
\documentclass[]{../../../../tex/classes/styledArticle}
\usepackage{../../../../tex/packages/hebrewSupport}
\usepackage{../../../../tex/packages/mathShortcuts}
\usepackage{../../../../tex/packages/theoremsSupport}
\usepackage{../../../../tex/packages/enfancysections}

%! ~~~ Document ~~~

\author{שחר פרץ}
\title{\textit{לינארית 2א 15}}
\date{9 ביוני 2025}
\begin{document}
    \maketitle
    \section{\en{Linear Transformation About Inner Product Vector Spaces}}
    \defi{$V$ ממ''פ ו־$T \co V \to V$ ט,ל. אז $T$ נקרתא \textit{סימטרית} ($\F = \R$) או \textit{הרמטית} ($\F = \C$) אם $\forall u, v \in V\la Tu \mid v \ra = \la u, Tv \ra$
    באופן כללי, העתקה כזו תקרא \textit{צמודה לעצמה}. }
    
    \textbf{דוגמה. }עבור $V = \R^{n}$ ו־$\smut $ מ''פ סטנדרטית, עבור $A \in M_n(\R)$ מתקיים $T_A \co V \to V$ ט''ל. נקרא מתי היא מודה לעצמה – $\mut{v}{u} = v^{T}u$: 
    \[ \mut{T_Av}{u} = (Av)^{T}u = v^{T}A^Tu = \mut{v}{A^{T}u} \]
    ז''א אם $A = A^{T}$ אז $T_A$ סימטרית, כלומר $A$ מטריצה סימטרית. גם הכיוון השני נכון: אם $T \co V \to V$ סימטרית אז ע''י בחירת בסיס נקבל $[T]^{B}_B$ גם היא סימטרית. 
    \theo{העתקה סימטרית אמ''מ היא דומה למטריצה סימטרית. }
    
    \theo{יהיו $T, S \co V \to V$ צמודות לעצמן. אז: 
    \begin{enumerate}
        \item $\ag T, T + S$ צמודות לעצמן. 
        \item המכפלה $S\circ T$ צמודה לעצמה אמ''מ $ST = TS$
        \item אם $p$ פולינום מעל $\F$ אז $p(T)$ צמדוה לעצמה. 
    \end{enumerate}}
קל לראות ש־$1 + 2 \implies 3$. $1$ נובע ישירות מהגדרה. נוכיח את $2$. \begin{proof}[הוכחה ל־2.]
    נניח $S \circ T$ צמודה לעצמה. בהנחות המשפט ידוע $S, T$ צמודות לעצמן. נקבל: 
        \[ \mut{(S \circ T)v}{u} = \mut{v}{STu} = \mut{Sv}{Tu} = \mut{TSv}{u} \implies \mut{(ST - TS)v}{u} = 0 \quad \forall v, u \]
        נסיק: 
        \[ \implies \forall v\mut{(ST - TS)v}{(ST - TS)v} = 0  \implies (ST - TS)v = 0 \implies STv = TSv \implies \top \]
        מהכיוון השני: 
        \[ \mut{STv}{u} = \mut{S(Tv)}{u} = \mut{v}{TSu} = \mut{v}{STu} \]
    \end{proof}
    \defi{$T \co V \to V$ תקרא חיובית/אי־שלילית/שלילית/אי־חיובית אם: 
    \begin{itemize}
        \item חיובית: \hfil $\mut{Tv}{v} > 0$
        \item שלילית: וכו'
    \end{itemize}}
    \theo{אם $T$ חיובית, אז היא הפיכה (כנ''ל לשלילית)} \begin{proof}
        נניח ש־$T$ לא הפיכה, נקרא שהיא לא חיובית. קיים $0 \neq v \in V$. אז $v \in \ker T$, אז $\mut{Tv}{v} = \mut{0}{v} = 0$, בסתירה לכך ש־$T$ חיובית. 
    \end{proof}
    \theo{נניח ש־$S$ צמודה לעצמה, אז $S^{2}$ צמודה לעצמה ואי־שלילית. }
    \begin{proof}
        ממשפט קודם $S^{2}$ צמודה לעצמה. נוכיח אי־שלילית: 
        \[ \forall 0 \neq v \in V \co \mut{S^{2}v}{v} = \mut{Sv}{Sv} = \norm{Sv}^{2} \ge 0 \]
    \end{proof}
    \defi{פולינום $p \in \R[x]$ יקרא חיובי אם $\forall x \in \R p(x) > 0$. }
    \textbf{מסקנה. }נניח $p(x) \in \R[x]$ חיובי, ו־$T \co V \to V$ צמודה לעצמה, אז $p(T)$ חיובית גם־כן, וצמודה לעצמה. 
    \lem{אם $p \in \R[x]$ חיובי, אז קיימים $g_1 \dots g_k \in \R[x]$ וכן $0 < c \in \R$ כך ש־$p(x) = \sum_{i = 1}^{k}g^2_i(x) + c$}
    רעיון להוכחת הלמה: מעל $\C$ זה מתפרק, ונוכל לכתוב $p(x) = a_n\prod_{j = 1}^{s} (x - i\ag_j)(x + i\ag_j)$ (מעל $\R$ כל פולינום מתפרק לגורמים ריבועיים, ואם כל שורשיו מרוכבים, כל גורמיו ריבועיים. ). הרעיון הוא להוכיח את הטענה ש־$g^2 h \bar h = g_1^2 + g_2^2$.
    \begin{proof}[הוכחה (של המשפט, לא של הלמה). ]
        יהי $0 \neq v \in V$. אז: 
        \[ \mut{p(T)v}{v} = \underbrace{\mut{\sum_{i = 1}^{k}g_i^2(T)v}{v}}_{\mathclap{\sum_{i = 1}^{k} \mut{g_i^2(T)v}{v}\ge 0}} + \overbrace{c\mut{v}{v}}^{c\norm{v}^2 > 0} \ge 0 \]
    \end{proof}
    \textbf{מסקנה. }אם $T \co V \to V$ צמודה לעצמה ו־$p(x) \in \R[x]$ פולינום חיובי, אז $p(T)$ הפיכה. 
    \begin{proof}
        ``תסתכלו על צד ימין של הלוח'' $\sim$ המרצה
    \end{proof}
    \theo{נניח ש־$T \co V \to V$ סימטרית (צמודה מעצה מעל $\R$/המייצגת סימטרית) ויהי $m_T(x)$ הפולינום המינימלי של $T$. אז $m_T$ מתפרק לגורמים לינארים. בנוסף, הם שונים זה מזה. }
    \textbf{מסקנה. }$T$ סימטרית לכסינה. 
    \begin{proof}
        נניח בשלילה קיום $p \mid m_T$ ו־$\deg p \ge 2$, $p $אי־פריק. בה''כ נניח ש־$p$ חיובי (אין לו שורש ב־$\R$, לכן נמצא כולו מעל/מתחת לציר ה־$x$). אז אפשר לכתוב את $m_T$ כ־$m_T = p \cdot g$ כלשהו. ידוע $p(T) \neq 0$ כי $m_T$ מינימלי מדרגה גבוהה יותר. 
        אזי: 
        \[ 0 = m_T(T) = \underbrace{p(T)}_{\neq 0} \cdot g(T) \implies g(T) = 0 \]
        בסתירה למינימליות של $m_T$. סה''כ $m_T$ אכן מתפרק לגורמים לינארים. עתה יש להראות שהגורמים הלינארים שלו זרים. נניח ש־$T$ סימטרית. ניעזר בלמה המופיע מיד אחרי ההוכחה הזו. נניח בשלילה שהם לא כולם שונים, אז $m_T(x) = (x - \lg)^2g(x)$ ואז: 
        \[ 0 = m_T(T)v = (T - \lg I)^{2}g(T) \implies \omega = g(T)v, \ (T - \lg I)^2\omega = 0 \]
        לכן בפרט $(T - \lg I)\omega = 0$ מהסעיף הקודם. סה''כ $\forall v \in V \co (T - \lg I)g(T) = 0$ וסתירה למינימליות. 
    \end{proof}
    \lem{נניח $T$ סממטרית ו־$\lg \in \R$, אם $(T - \lg I)^{2} = 0$ אז $T - \lg I = 0$. }\begin{proof}
        ידוע: 
        \[ \forall v \co 0 = \mut{(T - \lg I)^{2}v}{v} = \mut{(T - \lg I)v}{(T - \lg I)v} = \norm{(T - \lg I)v}^2 \implies (T - \lg I)v = 0 \]
    \end{proof}

    
    
    
    
    \ndoc
\end{document}
